\chapter{Resultados Experimentales}

\section{Validación dimensional independiente de objetos patrón}

Como paso previo a la evaluación del sistema CMOS, se realizó la caracterización dimensional independiente de las microesferas de poliestireno de diámetro nominal $2\,\mu m$, empleando microscopía óptica calibrada y microscopía electrónica de barrido (SEM).

Las imágenes obtenidas mediante SEM permitieron observar con alta definición la morfología esférica de las partículas y realizar mediciones directas del diámetro aparente.

\begin{figure}[H]
    \centering
    \includegraphics[width=0.5\textwidth]{fig/Fig_05_01_SEM_Microesferas.png}
    \caption{Micrografía SEM representativa de microesferas de poliestireno (2 µm nominal).}
    \label{fig:sem_microesferas}
\end{figure}

A partir de la medición de múltiples partículas, se obtuvo un diámetro promedio consistente con el valor nominal dentro del margen de dispersión esperado.

\begin{table}[H]
\centering
\caption{Resultados de medición dimensional mediante SEM}
\begin{tabular}{|c|c|}
\hline
\textbf{Número de muestras} & $n$ \\
Diámetro promedio & $\bar{d}_{SEM}$ \\
Desviación estándar & $\sigma_{SEM}$ \\
\hline
\end{tabular}
\end{table}

Estos resultados constituyen la referencia dimensional independiente utilizada para contrastar las mediciones obtenidas mediante el sistema CMOS.
\section{Adquisición directa mediante sensor CMOS en configuración lensless}

Se realizaron capturas directas de las microesferas empleando el sistema CMOS OV5647 en configuración de microscopía sin lentes.

Las imágenes adquiridas muestran estructuras circulares de bajo contraste, cercanas al límite de resolución impuesto por el tamaño de píxel ($\sim 1.4\,\mu m$).

\begin{figure}[H]
    \centering
    \includegraphics[width=0.5\textwidth]{fig/Fig_05_02_CMOS_RAW.png}
    \caption{Imagen RAW adquirida con el sensor CMOS en configuración lensless.}
    \label{fig:cmos_raw}
\end{figure}

Se observa que las microesferas no presentan bordes perfectamente definidos, sino patrones suavizados influenciados por fenómenos de difracción y por la respuesta espacial del sensor.
\section{Mejora perceptual mediante procesamiento computacional}

Con el objetivo de evaluar el impacto del procesamiento digital sobre la interpretabilidad de la imagen, se aplicaron técnicas de reconstrucción y super-resolución basadas en aprendizaje profundo.

\subsection{Reconstrucción FPM}

Se realizó una reconstrucción preliminar utilizando múltiples capturas bajo variación angular de iluminación.

\begin{figure}[H]
    \centering
    \includegraphics[width=0.5\textwidth]{fig/Fig_05_03_FPM.png}
    \caption{Resultado de reconstrucción mediante FPM.}
    \label{fig:fpm_result}
\end{figure}

Se observa mejora en contraste y definición visual, aunque la resolución física continúa limitada por el tamaño de píxel.

\subsection{Escalamiento mediante RealESRGAN}

Se aplicó super-resolución digital sobre la imagen base.

\begin{figure}[H]
    \centering
    \includegraphics[width=0.85\textwidth]{fig/Fig_05_04_SuperResolucion.png}
    \caption{Comparación entre imagen original y resultado tras super-resolución digital.}
    \label{fig:sr_result}
\end{figure}

Si bien la percepción visual mejora, se reconoce que estas técnicas no incrementan la resolución física real del sistema.
\section{Estimación experimental de la escala espacial}

Se procedió a realizar mediciones digitales del diámetro aparente de microesferas seleccionadas en regiones de interés (ROI).

El procedimiento consistió en:

\begin{itemize}
    \item Selección de microesferas aisladas.
    \item Medición del diámetro en píxeles.
    \item Conversión a unidades físicas mediante una escala preliminar ($\mu m/pixel$).
\end{itemize}

\begin{figure}[H]
    \centering
    \includegraphics[width=0.55\textwidth]{fig/Fig_05_05_Medicion_Diametro_CMOS.png}
    \caption{Ejemplo de medición digital del diámetro aparente en píxeles.}
    \label{fig:diametro_cmos}
\end{figure}

A partir de la comparación con el diámetro promedio obtenido por SEM, se estimó una escala espacial aproximada:

\[
\text{Escala experimental} \approx 0.17 \,\mu m/pixel
\]

La diferencia relativa entre mediciones se mantuvo dentro de un margen coherente considerando las limitaciones físicas del sistema.
\section{Análisis de limitaciones físicas y computacionales}

Los resultados experimentales evidencian que el sistema es capaz de detectar estructuras del orden de $2\,\mu m$, aunque operando cerca del límite de muestreo impuesto por el tamaño de píxel.

Las principales limitaciones identificadas incluyen:

\begin{itemize}
    \item Influencia significativa de la difracción en la definición de bordes.
    \item Dependencia crítica de la distancia muestra–sensor.
    \item Sensibilidad a variaciones en iluminación.
    \item Resolución efectiva limitada por el tamaño de píxel.
\end{itemize}

Estos factores explican las discrepancias observadas entre la imagen CMOS y las micrografías SEM, y refuerzan la necesidad de complementar la validación con patrones litográficos de mayor robustez geométrica.
\section{Síntesis de resultados y cumplimiento de objetivos}

Los resultados obtenidos permiten afirmar que:

\begin{itemize}
    \item Se implementó exitosamente un sistema de adquisición basado en sensor CMOS.
    \item Se estableció una metodología de validación metrológica preliminar.
    \item Se contrastaron mediciones con técnicas independientes (óptica y SEM).
    \item Se identificaron límites físicos asociados a resolución y difracción.
\end{itemize}

En consecuencia, el sistema demuestra coherencia dimensional preliminar y establece una base técnica sólida para futuras etapas de optimización del payload.
